% id: 91-05
% book: 베이직쎈 미적분1 · 05 도함수의 활용 (2)
% section: 학교 시험 기출로 실전 감각 UP!
% figure: tikz:fig-91-05
% answer: 
% answer_source: 
% variant_level: 0 (원본 전사)
% 이 파일은 build.mjs 가 items.json 에서 만든다. 고칠 때는 items.json 을 고친다.
\smash{\raisebox{1.5mm}{\dmkichul{베이직쎈 미적분1 91쪽 05번 · 꼭 나와요}}}
구간 $(-5{,}\mkern3mu\mathopen{}6)$에서 정의된 함수 $f(x)$의 도함수 $y=f'(x)$의 그래프가 다음 그림과 같을 때, 옳은 것만을 \textbf{보기}에서 있는 대로 고른 것은? 
\par\smallskip\begin{center}\resizebox{\ifdim\width>\linewidth\linewidth\else\width\fi}{!}{% 91-05(91쪽) — 구간 $(-5,\,6)$ 에서 정의된 $f(x)$ 의 도함수 $y=f'(x)$ 의 그래프($x=-4$ 에서 $x$축을 지나 내려가 $x=-2$ 에서 극소, $x=1$ 에서 $x$축에 닿았다가 다시 내려가 $x=3$ 에서 극소, $x=5$ 에서 $x$축을 지나 올라감). 스캔 화질이 낮아 TikZ 로 다시 그림.
% 원문에 있는 것만 옮긴다: 곡선 · 양 끝의 빈 점 ○(구간이 열려 있음) · $x=-5$, $-2$, 3, 6 의 안내 점선 · $-4$ 를 가리키는 작은 화살표 · 라벨 $-5$, $-4$, $-2$, O, 1, 3, 5, 6, x, y, y=f'(x). $x=1$ 에서 곡선이 $x$축에 닿기만 하는 것(부호가 바뀌지 않음)을 원문 그대로 옮겼다.
\begin{tikzpicture}[x=0.52cm, y=0.52cm, line width=0.6pt]
  \draw[->, >=stealth, line width=0.7pt] (-6.2,0) -- (6.9,0) node[below=1pt, font=\small] {$x$};
  \draw[->, >=stealth, line width=0.7pt] (0,-3.1) -- (0,3.3) node[left=1pt, font=\small] {$y$};
  \draw[densely dotted, line width=0.4pt] (-5,0) -- (-5,1.38);
  \draw[densely dotted, line width=0.4pt] (-2,0) -- (-2,-2.62);
  \draw[densely dotted, line width=0.4pt] (3,0) -- (3,-1.81);
  \draw[densely dotted, line width=0.4pt] (6,0) -- (6,1.34);
  \draw plot[smooth, tension=0.75] coordinates {(-5,1.38) (-4.6,0.72) (-4,0) (-3.6,-0.72) (-3.2,-1.42) (-2.8,-2.02) (-2.4,-2.44) (-2,-2.62) (-1.6,-2.55) (-1.2,-2.3) (-0.8,-1.92) (-0.4,-1.5) (0,-1.06) (0.4,-0.6) (0.7,-0.22) (1,0) (1.3,-0.2) (1.6,-0.55) (2,-1.05) (2.4,-1.5) (2.7,-1.72) (3,-1.81) (3.4,-1.7) (3.8,-1.4) (4.2,-0.95) (4.6,-0.5) (5,0) (5.4,0.56) (5.7,1) (6,1.34)};
  \draw[fill=white, line width=0.5pt] (-5,1.38) circle (1.7pt);
  \draw[fill=white, line width=0.5pt] (6,1.34) circle (1.7pt);
  \draw[->, >=stealth, line width=0.4pt] (-4.5,-1.15) to[out=75, in=250] (-4.25,-0.22);
  \node[font=\small] at (-4.62,-1.62) {$-4$};
  \node[below left=-1pt and -5pt, font=\small] at (-5,0) {$-5$};
  \node[above=1pt, font=\small] at (-2,0) {$-2$};
  \node[below left=-1pt and -4pt, font=\small] at (0,0) {$\pt{O}$};
  \node[above=2pt, font=\small] at (1,0) {$1$};
  \node[above=1pt, font=\small] at (3,0) {$3$};
  \node[below=1pt, font=\small] at (5,0) {$5$};
  \node[below=1pt, font=\small] at (6.05,0) {$6$};
  \node[font=\small] at (5.6,2.05) {$y=f'(x)$};
\end{tikzpicture}
}\end{center}
\noindent \bogi{ㄱ. $f(x)$는 $x=1$에서 극댓값을 갖는다.\\ ㄴ. $f(x)$는 $x=5$에서 극솟값을 갖는다.\\ ㄷ. $f(x)$가 극값을 갖는 점은 3개이다.}
\dmchoicesauto{}{ㄱ}{ㄴ}{ㄷ}{ㄱ, ㄴ}{ㄴ, ㄷ}
